\begin{longtable}{TD}
\caption{Task entity terms.}
\label{tab:terms-task-entities}\\

\toprule
Term & Definition \\
\midrule
\endfirsthead

\toprule
Term & Definition \\
\midrule
\endhead

\bottomrule
\endfoot

\termrow{task}{
A generic unit of work handled by a queue-worker system. The term is used only
when the distinction between task kind, task instance, and task attempt is not
important.
}

\termrow{task kind}{
The semantic type of work submitted to a queue, such as copying a replication
partition, rebuilding a cache entry, resizing an image, repairing a database
record, dispatching a webhook, or validating a document. A task kind describes
what work is requested, but it does not fully determine how that work behaves
under every execution context.
}

\termrow{task instance}{
A concrete logical unit of work placed into a queue. A task instance belongs to
a task kind and may produce one or more execution attempts.
}

\termrow{task attempt}{
One concrete execution attempt of a task instance. Retries, lease expiration,
worker failure, cancellation, or rescheduling may cause a single task instance
to have multiple attempts. Resource demand, worker occupancy, observations,
outcomes, and residuals are usually measured or estimated at attempt
granularity.
}

\termrow{task family}{
A group of task kinds that share a similar behavior-graph structure or profiling
model. A task family can be used when one graph template applies to several
related task kinds.
}

\termrow{task identity}{
The identifying metadata of a task, including task kind, task version, queue,
producer, tenant, priority, and deadline. Task identity identifies the work but
does not itself define the cost of executing it.
}

\termrow{task version}{
The version of a task kind, handler contract, payload schema, or behavior
manifest. Task versions are useful for compatibility between historical
profiles, manifests, and schema versions.
}

\termrow{queue}{
The queue through which a task instance is submitted, stored, reserved, and
eventually consumed by a worker. A queue name may provide useful context, but it
is not a complete behavior model.
}

\termrow{producer}{
The component that creates and submits a task instance to a queue. The producer
may provide metadata or declared hints, but it does not fully determine the
execution context.
}

\termrow{worker}{
An executor that dequeues, reserves, executes, retries, acknowledges, discards,
or fails task attempts. In this paper, a worker is treated as an execution
boundary rather than merely a message consumer.
}

\termrow{worker pool}{
A set of workers that execute task attempts. Worker-pool capacity and
worker-pool occupancy are central to queue-worker profiling.
}

\termrow{tenant}{
A user, customer, namespace, application, or logical group that consumes
execution capacity. Tenants are relevant for isolation, fairness, consequence
analysis, and operational risk.
}

\termrow{tenant class}{
A category of tenants based on SLO, priority, plan, criticality, or operational
risk. A tenant class is distinct from a task kind.
}

\termrow{priority}{
A scheduling or ordering attribute assigned to a task instance. Priority may
influence profiling or admission decisions, but it is not itself resource demand
or operational risk.
}

\termrow{deadline}{
A completion target or time constraint associated with a task instance. A
deadline may affect consequence and risk, but it is distinct from observed
latency.
}

\end{longtable}
